\documentclass[aspectratio=169]{beamer}

% Shared preamble for Claude Code slide decks
% Usage: % Shared preamble for Claude Code slide decks
% Usage: % Shared preamble for Claude Code slide decks
% Usage: \input{preamble} after \documentclass[aspectratio=169]{beamer}

\usetheme{default}
\usecolortheme{dove}

\setbeamertemplate{navigation symbols}{}
\setbeamertemplate{footline}{%
  \hfill{\large\insertframenumber\,/\,\inserttotalframenumber}\hspace{0.8em}\vspace{0.5em}%
}

\definecolor{popblue}{RGB}{52, 101, 164}
\definecolor{sampred}{RGB}{204, 0, 0}
\definecolor{paramgreen}{RGB}{0, 140, 70}
\definecolor{lightbg}{RGB}{245, 245, 250}
\definecolor{warnred}{RGB}{180, 40, 40}
\definecolor{orange1}{RGB}{220, 120, 0}
\definecolor{violet1}{RGB}{120, 50, 160}

\setbeamercolor{frametitle}{fg=popblue}
\setbeamercolor{title}{fg=popblue}
\setbeamercolor{block title}{fg=white,bg=popblue}
\setbeamercolor{block body}{bg=lightbg}
\setbeamercolor{block title alerted}{fg=white,bg=warnred}
\setbeamercolor{block body alerted}{bg=lightbg}
\setbeamercolor{block title example}{fg=white,bg=paramgreen}
\setbeamercolor{block body example}{bg=lightbg}

\usepackage{listings}
\usepackage[T1]{fontenc}
\usepackage{hyperref}
\usepackage{tikz}
\usepackage{booktabs}
\usetikzlibrary{shapes, arrows.meta, positioning, calc}

\lstset{
  basicstyle=\ttfamily\scriptsize,
  keywordstyle=\color{popblue}\bfseries,
  stringstyle=\color{paramgreen},
  commentstyle=\color{gray},
  breaklines=true,
  frame=single,
  backgroundcolor=\color{lightbg},
  rulecolor=\color{popblue!30},
  showstringspaces=false,
  columns=fullflexible,
}
 after \documentclass[aspectratio=169]{beamer}

\usetheme{default}
\usecolortheme{dove}

\setbeamertemplate{navigation symbols}{}
\setbeamertemplate{footline}{%
  \hfill{\large\insertframenumber\,/\,\inserttotalframenumber}\hspace{0.8em}\vspace{0.5em}%
}

\definecolor{popblue}{RGB}{52, 101, 164}
\definecolor{sampred}{RGB}{204, 0, 0}
\definecolor{paramgreen}{RGB}{0, 140, 70}
\definecolor{lightbg}{RGB}{245, 245, 250}
\definecolor{warnred}{RGB}{180, 40, 40}
\definecolor{orange1}{RGB}{220, 120, 0}
\definecolor{violet1}{RGB}{120, 50, 160}

\setbeamercolor{frametitle}{fg=popblue}
\setbeamercolor{title}{fg=popblue}
\setbeamercolor{block title}{fg=white,bg=popblue}
\setbeamercolor{block body}{bg=lightbg}
\setbeamercolor{block title alerted}{fg=white,bg=warnred}
\setbeamercolor{block body alerted}{bg=lightbg}
\setbeamercolor{block title example}{fg=white,bg=paramgreen}
\setbeamercolor{block body example}{bg=lightbg}

\usepackage{listings}
\usepackage[T1]{fontenc}
\usepackage{hyperref}
\usepackage{tikz}
\usepackage{booktabs}
\usetikzlibrary{shapes, arrows.meta, positioning, calc}

\lstset{
  basicstyle=\ttfamily\scriptsize,
  keywordstyle=\color{popblue}\bfseries,
  stringstyle=\color{paramgreen},
  commentstyle=\color{gray},
  breaklines=true,
  frame=single,
  backgroundcolor=\color{lightbg},
  rulecolor=\color{popblue!30},
  showstringspaces=false,
  columns=fullflexible,
}
 after \documentclass[aspectratio=169]{beamer}

\usetheme{default}
\usecolortheme{dove}

\setbeamertemplate{navigation symbols}{}
\setbeamertemplate{footline}{%
  \hfill{\large\insertframenumber\,/\,\inserttotalframenumber}\hspace{0.8em}\vspace{0.5em}%
}

\definecolor{popblue}{RGB}{52, 101, 164}
\definecolor{sampred}{RGB}{204, 0, 0}
\definecolor{paramgreen}{RGB}{0, 140, 70}
\definecolor{lightbg}{RGB}{245, 245, 250}
\definecolor{warnred}{RGB}{180, 40, 40}
\definecolor{orange1}{RGB}{220, 120, 0}
\definecolor{violet1}{RGB}{120, 50, 160}

\setbeamercolor{frametitle}{fg=popblue}
\setbeamercolor{title}{fg=popblue}
\setbeamercolor{block title}{fg=white,bg=popblue}
\setbeamercolor{block body}{bg=lightbg}
\setbeamercolor{block title alerted}{fg=white,bg=warnred}
\setbeamercolor{block body alerted}{bg=lightbg}
\setbeamercolor{block title example}{fg=white,bg=paramgreen}
\setbeamercolor{block body example}{bg=lightbg}

\usepackage{listings}
\usepackage[T1]{fontenc}
\usepackage{hyperref}
\usepackage{tikz}
\usepackage{booktabs}
\usetikzlibrary{shapes, arrows.meta, positioning, calc}

\lstset{
  basicstyle=\ttfamily\scriptsize,
  keywordstyle=\color{popblue}\bfseries,
  stringstyle=\color{paramgreen},
  commentstyle=\color{gray},
  breaklines=true,
  frame=single,
  backgroundcolor=\color{lightbg},
  rulecolor=\color{popblue!30},
  showstringspaces=false,
  columns=fullflexible,
}


\title{Claude 101}
\subtitle{Prompting $\cdot$ Projects $\cdot$ Artifacts $\cdot$ Skills $\cdot$ Research}
\date{}

\begin{document}

\begin{frame}
\titlepage
\end{frame}

% ──────────────────────────────────────────────
\begin{frame}{Course Overview}
\begin{columns}[T]
\column{0.33\textwidth}
\textbf{\color{popblue}Foundations}
\begin{itemize}
  \item What is Claude?
  \item First conversation
  \item Getting better results
  \item Desktop app modes
\end{itemize}

\column{0.33\textwidth}
\textbf{\color{paramgreen}Features}
\begin{itemize}
  \item Projects
  \item Artifacts
  \item Skills
  \item Connectors \& MCP
  \item Enterprise search
  \item Research mode
\end{itemize}

\column{0.33\textwidth}
\textbf{\color{violet1}Application}
\begin{itemize}
  \item Use cases by role
  \item Other ways to work\\(Code, Slack, Excel)
\end{itemize}
\end{columns}
\end{frame}

% ──────────────────────────────────────────────
\begin{frame}{What Is Claude?}
\textbf{Claude} = AI assistant by Anthropic, designed to be \textbf{helpful, harmless, and honest}.

\bigskip
\textbf{Built with Constitutional AI:}
\begin{itemize}
  \item Trained on a set of principles (a ``constitution'')
  \item Self-evaluates responses against those principles
  \item Aims for safe, balanced, truthful outputs
\end{itemize}

\bigskip
\textbf{Key capabilities:}
\begin{itemize}
  \item Long-context understanding (200K--1M tokens depending on model)
  \item Strong reasoning and analysis
  \item Code generation and review
  \item Document creation and summarization
  \item Tool use and integration
  \item Adaptive thinking --- adjusts reasoning depth to task complexity
\end{itemize}

\bigskip
\textbf{Access:} claude.ai, desktop app, API, Claude Code CLI, Slack, Excel\ldots
\end{frame}

% ──────────────────────────────────────────────
\begin{frame}{Your First Conversation: Prompting}
\textbf{The 3-part prompt structure:}

\begin{enumerate}
  \item \textbf{\color{popblue}Set the stage} --- give Claude a role and context
  \item \textbf{\color{paramgreen}Define the task} --- be specific about what you need
  \item \textbf{\color{orange1}Specify rules} --- format, constraints, what to avoid
\end{enumerate}

\bigskip
\begin{exampleblock}{Example}
``You are a senior data analyst. Analyze this CSV of quarterly sales data. Create a summary table showing trends by region. Use bullet points, not paragraphs. Flag any anomalies.''
\end{exampleblock}

\bigskip
\textbf{Uploads:} paste images, attach files (PDFs, CSVs, code) directly in the conversation.

\textbf{Iteration:} Claude's first answer is a starting point --- refine and redirect.
\end{frame}

% ──────────────────────────────────────────────
\begin{frame}{Getting Better Results}
\begin{columns}[T]
\column{0.5\textwidth}
\textbf{\color{warnred}Common challenges:}
\begin{itemize}
  \item Too generic / vague output
  \item Wrong format or structure
  \item Missing key details
  \item Hallucinated information
  \item Over-verbose responses
\end{itemize}

\bigskip
\textbf{\color{popblue}AI Fluency 4D Framework:}
\begin{enumerate}
  \item \textbf{Discover} --- explore what's possible
  \item \textbf{Define} --- scope the task precisely
  \item \textbf{Develop} --- iterate on outputs
  \item \textbf{Deliver} --- finalize and apply
\end{enumerate}

\column{0.5\textwidth}
\textbf{\color{paramgreen}Tips for better prompts:}
\begin{itemize}
  \item Be specific, not vague
  \item Provide examples of desired output
  \item Break complex tasks into steps
  \item Specify format and length
  \item Give context about your audience
  \item Use \textbf{evals} to measure quality systematically
\end{itemize}
\end{columns}
\end{frame}

% ──────────────────────────────────────────────
\begin{frame}{Desktop App: Chat, Cowork, Code}
\begin{table}
\centering\small
\begin{tabular}{lccc}
 & \textbf{Chat} & \textbf{Cowork} & \textbf{Code} \\
\hline
Interaction & Conversational & Delegation & CLI/IDE \\
Best for & Q\&A, drafting & Complex tasks & Software dev \\
Autonomy & Low & High & High \\
File access & Upload & Folder access & Full repo \\
Subagents & No & Yes & Yes \\
\end{tabular}
\end{table}

\bigskip
\textbf{Quick entry:} launch from system tray, dictation support.

\begin{block}{Choose the right mode}
\textbf{Chat} for quick questions. \textbf{Cowork} for ``do this for me'' tasks. \textbf{Code} for development.
\end{block}
\end{frame}

% ──────────────────────────────────────────────
\begin{frame}{Projects}
\textbf{Projects} = persistent knowledge bases for Claude.

\bigskip
\begin{columns}[T]
\column{0.5\textwidth}
\textbf{What you can add:}
\begin{itemize}
  \item Reference documents (PDFs, docs)
  \item Custom instructions
  \item Style guides and templates
  \item Data files
\end{itemize}

\bigskip
\textbf{How it works:}
\begin{itemize}
  \item RAG (Retrieval Augmented Generation)
  \item Claude searches project files for relevant context
  \item No need to re-upload every conversation
\end{itemize}

\column{0.5\textwidth}
\textbf{Collaboration:}
\begin{itemize}
  \item Share projects with team members
  \item Everyone gets the same context
  \item Consistent instructions across users
\end{itemize}

\bigskip
\begin{exampleblock}{Use cases}
\begin{itemize}
  \item Company style guide project
  \item Product documentation hub
  \item Research paper collection
  \item Client-specific knowledge base
\end{itemize}
\end{exampleblock}
\end{columns}
\end{frame}

% ──────────────────────────────────────────────
\begin{frame}{Artifacts}
\textbf{Artifacts} = rich, interactive outputs Claude creates alongside responses.

\bigskip
\textbf{Types of artifacts:}
\begin{columns}[T]
\column{0.5\textwidth}
\begin{itemize}
  \item \textbf{Documents} --- formatted text, reports
  \item \textbf{Code} --- syntax-highlighted, runnable
  \item \textbf{HTML/CSS} --- rendered web pages
  \item \textbf{SVG} --- vector graphics
\end{itemize}

\column{0.5\textwidth}
\begin{itemize}
  \item \textbf{Mermaid} --- diagrams and flowcharts
  \item \textbf{React components} --- interactive UIs
  \item \textbf{PowerPoint} --- downloadable \texttt{.pptx} files
  \item Shareable, publishable, MCP-connected
\end{itemize}
\end{columns}

\bigskip
\begin{block}{Workflow}
Ask Claude to create $\to$ preview in the artifact panel $\to$ iterate $\to$ share or publish.
\end{block}
\end{frame}

% ──────────────────────────────────────────────
\begin{frame}{Skills \& Connectors}
\begin{columns}[T]
\column{0.5\textwidth}
\textbf{\color{popblue}Skills}
\begin{itemize}
  \item \textbf{Anthropic skills} --- built-in capabilities
  \item \textbf{Custom skills} --- user/team-defined (see \emph{Agent Skills} deck)
  \item Enable from settings
  \item Activate automatically based on context
\end{itemize}

\bigskip
\textbf{Skills vs.\ Projects:}\\
Skills = \emph{how} to do things\\
Projects = \emph{what} to know

\column{0.5\textwidth}
\textbf{\color{paramgreen}Connectors}
\begin{itemize}
  \item Google Drive, Notion, Slack\ldots
  \item MCP (Model Context Protocol)
  \item Web connectors for any service
  \item Desktop extensions
\end{itemize}

\bigskip
\textbf{Enterprise Search:}
\begin{itemize}
  \item Admin configures data sources
  \item Users query across all connected systems
  \item Security: respects original permissions
\end{itemize}
\end{columns}
\end{frame}

% ──────────────────────────────────────────────
\begin{frame}{Research Mode}
\textbf{Research mode} = multi-step investigation with extended thinking.

\bigskip
\textbf{How it works:}
\begin{enumerate}
  \item Claude breaks your question into sub-questions
  \item Searches across available sources
  \item Synthesizes findings with \textbf{citations}
  \item Produces a structured research report
\end{enumerate}

\bigskip
\textbf{Best for:}
\begin{itemize}
  \item Complex questions requiring multiple sources (up to 45 min)
  \item Competitive analysis and market research
  \item Literature review and synthesis
  \item Due diligence and fact-checking
\end{itemize}

\begin{alertblock}{Tip}
Research mode uses extended thinking --- takes longer but produces deeper, more accurate results.
\end{alertblock}
\end{frame}

% ──────────────────────────────────────────────
\begin{frame}{Use Cases by Role}
\begin{columns}[T]
\column{0.5\textwidth}
\begin{itemize}
  \item \textbf{\color{popblue}Sales} --- prospect research, email drafting, proposal generation
  \item \textbf{\color{paramgreen}Marketing} --- content creation, campaign analysis, copywriting
  \item \textbf{\color{orange1}Finance} --- data analysis, report generation, modeling assistance
\end{itemize}

\column{0.5\textwidth}
\begin{itemize}
  \item \textbf{\color{violet1}HR} --- policy drafting, interview prep, training materials
  \item \textbf{\color{warnred}Legal} --- document review, contract analysis, compliance checks
  \item \textbf{\color{sampred}Research} --- literature synthesis, experiment design, data interpretation
\end{itemize}
\end{columns}

\bigskip
\textbf{Other interfaces:}
Claude Code (CLI) $\cdot$ Slack integration $\cdot$ Excel add-in $\cdot$ Chrome extension
\end{frame}

% ──────────────────────────────────────────────
\begin{frame}{Summary}
\begin{block}{Key Takeaways}
\begin{itemize}
  \item Claude: helpful, harmless, honest --- Constitutional AI, adaptive thinking
  \item Prompt = set stage + define task + specify rules; iterate
  \item Three modes: Chat (quick), Cowork (delegate), Code (dev)
  \item Projects = persistent knowledge; Artifacts = rich outputs
  \item Skills activate automatically; Connectors extend reach
  \item Research mode for deep multi-source investigation
  \item Available everywhere: web, desktop, CLI, API, Slack, Excel
\end{itemize}
\end{block}

\vfill
\centering
\Large\color{popblue} Questions?
\end{frame}

\end{document}
